%%============================ Compiler Directives =======================%%
%%                                                                        %%
% !TeX program = lualatex
% !TeX encoding = utf8
% !TeX spellcheck = uk_UA
%%                                                                        %%
%%============================== Клас документа ==========================%%
%%                                                                        %%
\documentclass[14pt]{extarticle}
\IfFileExists{ukrcorr.sty}{\usepackage{ukrcorr}}{}
%%                                                                        %%
%%========================== Мови, шрифти та кодування ===================%%
%%                                                                        %%
\usepackage{fontspec}
\setsansfont{CMU Sans Serif}%{Arial}
\setmainfont{CMU Serif}%{Times New Roman}
\setmonofont{CMU Typewriter Text}%{Consolas}
\defaultfontfeatures{Ligatures={TeX}}
\usepackage[math-style=TeX]{unicode-math}
\usepackage[english, russian, ukrainian]{babel}

%%                                                                        %%
%%============================= Геометрія сторінки =======================%%
%%                                                                        %%
\usepackage[%
	a4paper,%
	footskip=1cm,%
	headsep=0.3cm,%
	top=2cm, %поле сверху
	bottom=2cm, %поле снизу
	left=2cm, %поле ліворуч
	right=2cm, %поле праворуч
    ]{geometry}
%%                                                                        %%
%%============================== Інтерліньяж  ============================%%
%%                                                                        %%
\renewcommand{\baselinestretch}{1}
%-------------------------  Подавление висячих строк  --------------------%%
\clubpenalty =10000
\widowpenalty=10000
%---------------------------------Інтервали-------------------------------%%
\setlength{\parskip}{0.5ex}%
\setlength{\parindent}{2.5em}%
%%                                                                        %%
%%=========================== Математичні пакети і графіка ===============%%
%%                                                                        %%
\usepackage{amsmath}
\usepackage{graphicx}
\usepackage{floatflt}
%%                                                                        %%
%%========================== Гіперпосилення (href) =======================%%
%%                                                                        %%
\usepackage[%colorlinks=true,
	%urlcolor = blue, %Colour for external hyperlinks
	%linkcolor  = malina, %Colour of internal links
	%citecolor  = green, %Colour of citations
	bookmarks = true,
	bookmarksnumbered=true,
	unicode,
	linktoc = all,
	hypertexnames=false,
	pdftoolbar=false,
	pdfpagelayout=TwoPageRight,
	pdfauthor={Ponomarenko S.M. aka sergiokapone},
	pdfdisplaydoctitle=true,
	pdfencoding=auto
	]%
	{hyperref}
		\makeatletter
	\AtBeginDocument{
	\hypersetup{
		pdfinfo={
		Title={\@title},
		}
	}
	}
	\makeatother
%%                                                                        %%
%%============================ Заголовок та автори =======================%%
%%                                                                        %%
\title{}
\author{}
%%                                                                        %%
%%========================================================================%%


\begin{document}
\begin{align*}\label{}
    x &= R\sin\theta\cos\phi \\
    y &= R\sin\theta\sin\phi \\
    z &= R\cos\theta
\end{align*}

\begin{align*}\label{}
    dx &= R(\cos\theta\cos\phi d\theta - \sin\theta\sin\phi d\phi )\\
    dy &= R(\cos\theta\sin\phi d\theta + \sin\theta\cos\phi d\phi ) \\
    dz &= -R\sin\theta d\theta
\end{align*}

\begin{align*}\label{}
    dx^2 &= R^2( \cos^2\theta\cos^2\phi d\theta^2 +  \sin^2\theta\sin^2\phi d\phi^2 - 2\cos\theta\sin\theta\cos\phi\sin\phi d\theta d\phi )\\
    dy^2 &= R^2( \cos^2\theta\sin^2\phi d\theta^2 +  \sin^2\theta\cos^2\phi d\phi^2 + 2\cos\theta\sin\theta\cos\phi\sin\phi  d\theta d\phi)\\
\end{align*}

\begin{equation*}\label{}
    dx^2 + dy^2 = R^2(\cos^2\theta d\theta^2 + \sin^2\theta d\phi^2)
\end{equation*}

\begin{align*}\label{}
    xdx &= R^2(\sin\theta\cos\theta\cos^2\phi d\theta - \sin^2\theta\cos\phi\sin\phi d\phi) \\
    ydy &= R^2(\sin\theta\cos\theta\sin^2\phi d\theta + \sin^2\theta\cos\phi\sin\phi d\phi)
\end{align*}

\begin{equation*}\label{}
    xdx + ydy = R^2\sin\theta\cos\theta d\theta
\end{equation*}

\begin{equation*}\label{}
    dl^2 = R^2( d\theta^2 + \sin^2\theta d\phi^2 )
\end{equation*}

\begin{equation*}\label{}
    \theta = \chi
\end{equation*}

\begin{equation*}\label{}
    dl^2 = R^2( d\chi^2 + \sin^2\chi d\phi^2 )
\end{equation*}

\begin{equation*}\label{}
    \int dl = 2\pi R\sin\chi
\end{equation*}


\begin{equation*}\label{}
    \int dS = \int\limits_0^{2\pi} d\phi \int\limits_0^{\pi} R^2\sin\chi d\chi = 4\pi R^2
\end{equation*}

\begin{equation*}\label{}
    \int dS = \int\limits_0^{2\pi} d\phi \int\limits_0^{\chi} R^2\sin\chi d\chi = 2\pi R^2 (1- \cos\chi) = 2\pi R^2(1 - \cos\frac{r}{R})
\end{equation*}

\end{document}


